\chapter{TÌM KIẾM CỤC BỘ TRONG TUYẾN}
\label{chap:ls}

\section{Mục tiêu và vị trí trong vòng lặp}

Cải tiến \code{local\_search} tinh chỉnh từng con vừa sinh ra bằng các thay đổi nhỏ trong từng tuyến.
Nó chạy \emph{sau đột biến} và \emph{trước lần đánh giá chính thức} của con. Không phải con nào cũng
được tinh chỉnh: với mỗi con, bộ sinh ngẫu nhiên của NSGA-II bốc một số trong $[0,1)$, nhỏ hơn $p_{ls}$
thì chạy tìm kiếm cục bộ, ngược lại con giữ nguyên và được đánh giá như thuật toán gốc.

Tham số trong \code{config.yaml}:
\begin{Verbatim}
improvements:
  local_search: false
  local_search_params:
    p_ls: 0.2                 # xác suất chạy LS cho mỗi con
    max_passes: 2             # số lượt duyệt các tuyến tối đa mỗi lần gọi
    max_candidate_evals: 40   # trần cứng số ứng viên được chấm mỗi lần gọi
    max_accepted_moves: 5     # số bước được nhận tối đa mỗi lần gọi
    operators: [two_opt, or_opt]
    or_opt_lengths: [1, 2, 3]
    strategy: first           # first | best
    acceptance: weighted      # weighted | pareto
    weight_distribution: dirichlet   # dirichlet | equal
    improvement_epsilon: 1.0e-10
    split_drift_policy: reject       # reject | evaluate_actual
\end{Verbatim}

\section{Hai toán tử trong tuyến}

\subsection{2-opt}

\code{two\_opt(route, i, j)} với $0 \le i < j < |\text{route}|$ trả về bản sao của tuyến với đoạn từ vị
trí $i$ đến $j$ (gồm cả hai đầu) bị đảo ngược. Tuyến đầu vào không bị sửa.

\subsection{Or-opt}

\code{or\_opt(route, start, L, insert\_at)} cắt khối $L$ khách liên tiếp bắt đầu tại \code{start}, giữ
nguyên thứ tự trong khối, rồi chèn khối vào vị trí \code{insert\_at}. Vị trí chèn được tính trên mảng
\emph{đã bỏ khối}, nên luôn nằm trong $[0, |\text{route}| - L]$. Khi $\code{insert\_at} = \code{start}$,
kết quả trùng tuyến cũ (không đổi gì). Ví dụ trong đặc tả: tuyến $(1,2,3,4,5,6)$, khối $(2,3)$ có
$\code{start}=1$, $L=2$; bỏ khối còn $(1,4,5,6)$; chèn tại vị trí 3 được $(1,4,5,2,3,6)$, tức khối được
chuyển ra sau khách 5.

\begin{figure}[H]
\centering
\begin{tikzpicture}[
  box/.style={draw, minimum width=0.9cm, minimum height=0.75cm, font=\small},
  hi/.style={box, fill=blue1!20},
  ho/.style={box, fill=orange1!25}]
  % 2-opt
  \node[anchor=east, font=\small\bfseries] at (-0.3,0) {2-opt};
  \foreach \v/\s [count=\k from 0] in {3/box,7/hi,1/hi,9/hi,4/hi,6/box}
    \node[\s] (a\k) at (\k*0.9,0) {\v};
  \node[font=\footnotesize] at (2.25,0.75) {đảo đoạn $i=1$ đến $j=4$};
  \draw[-{Stealth}, thick] (5.4,0) -- (6.4,0);
  \foreach \v/\s [count=\k from 0] in {3/box,4/hi,9/hi,1/hi,7/hi,6/box}
    \node[\s] at (7+\k*0.9,0) {\v};
  % or-opt
  \node[anchor=east, font=\small\bfseries] at (-0.3,-1.8) {Or-opt};
  \foreach \v/\s [count=\k from 0] in {3/box,7/ho,1/ho,9/box,4/box,6/box}
    \node[\s] at (\k*0.9,-1.8) {\v};
  \node[font=\footnotesize] at (2.25,-1.05) {khối $(7,1)$: \code{start}=1, $L$=2};
  \draw[-{Stealth}, thick] (5.4,-1.8) -- (6.4,-1.8);
  \foreach \v/\s [count=\k from 0] in {3/box,9/box,4/box,7/ho,1/ho,6/box}
    \node[\s] at (7+\k*0.9,-1.8) {\v};
  \node[font=\footnotesize] at (9.7,-1.05) {chèn tại vị trí 3 của $(3,9,4,6)$};
\end{tikzpicture}
\caption{Hai toán tử trong tuyến trên tuyến $(3,7,1,9,4,6)$}
\label{fig:ops}
\end{figure}

\subsection{Cách bốc bước đi}

Mỗi bước đi được bốc ngẫu nhiên, không duyệt hết mọi cặp vị trí (duyệt hết có $O(m^2)$ bước 2-opt và
$O(m^2)$ bước Or-opt trên một tuyến dài $m$, nhân với số tuyến và số lần lặp sẽ quá đắt). Với tuyến dài
$m \ge 2$:
\begin{itemize}
  \item 2-opt: bốc $i$ đều trong $[0, m-2]$, rồi $j$ đều trong $[i+1, m-1]$;
  \item Or-opt: bốc $L$ trong các độ dài cho phép nhỏ hơn $m$, bốc \code{start} trong $[0, m-L]$, bốc
        vị trí chèn trong $m - L$ vị trí khác \code{start}. Cách bốc này loại sẵn trường hợp không đổi gì.
\end{itemize}
Tuyến chỉ có một khách không có bước đi nào và được bỏ qua.

\section{Quy trình một lần tìm kiếm cục bộ}

\begin{algorithm}[H]
\caption{Một lần tìm kiếm cục bộ trên hoán vị con $\pi$}
\label{alg:ls}
\KwIn{$\pi$, bộ sinh \textit{rng}, tham số, \textit{eval} (đúng bộ đánh giá NSGA-II dùng), \textit{meta}}
\KwOut{hoán vị $\pi^*$ và kết quả đánh giá chính xác của nó}
\If{tiêu chí là \texttt{weighted}}{$\mathbf{w} \gets$ \textit{rng}.dirichlet$(1,1,1)$ \tcp*{một lần cho cả lần gọi}}
$D \gets \textit{eval}(\pi, \textit{meta})$; $\pi^* \gets \pi$\;
\For{lượt $\gets 1$ \KwTo \texttt{max\_passes}}{
  $R \gets D$.routes; $E \gets$ các tuyến có ít nhất 2 khách; \lIf{$E = \emptyset$}{\textbf{dừng}}
  \ForEach{tuyến $r$ theo thứ tự \textit{rng}.permutation($E$)}{
    \ForEach{toán tử $o \in$ \texttt{operators}}{
      \lIf{hết ngân sách ứng viên hoặc đủ số bước được nhận}{\Return{$\pi^*, D$}}
      $r' \gets$ bốc một bước $o$ trên $R[r]$; $R' \gets R$ với $R[r]$ thay bằng $r'$\;
      $\pi' \gets$ nối các tuyến của $R'$; $D' \gets \textit{eval}(\pi', \textit{meta})$\;
      \If{$D'$.routes $\ne R'$ \textbf{and} chính sách là \texttt{reject}}{loại (lệch biên tuyến); \textbf{tiếp}}
      \lIf{$D'$ làm xấu ràng buộc số xe so với $D$}{loại; \textbf{tiếp}}
      \If{$D'$ tốt hơn $D$ theo tiêu chí}{
        $\pi^* \gets \pi'$; $D \gets D'$; cập nhật $R$\;
        \textbf{sang tuyến kế tiếp} \tcp*{first improvement}
      }
    }
  }
}
\Return{$\pi^*, D$}
\end{algorithm}

Ba điểm thiết kế quan trọng:
\begin{enumerate}
  \item \textbf{Bước đi luôn được đề xuất trên bộ tuyến của nghiệm hiện tại}, lấy từ kết quả bộ giải mã,
        không tự cắt lại. Mỗi bước chỉ đổi thứ tự trong \emph{một} tuyến, tập khách của tuyến đó giữ nguyên.
  \item \textbf{Mọi ứng viên được chấm lại bằng chính bộ đánh giá NSGA-II dùng}, trên toàn bộ hoán vị.
        Mô-đun không tự tính điểm riêng cho một tuyến rồi coi đó là độ thích nghi của cả cá thể.
  \item \textbf{Kết quả trả về là kết quả đánh giá thật} của hoán vị cuối. NSGA-II dùng lại kết quả này,
        không phải chấm lần nữa.
\end{enumerate}
Với chiến lược \code{best}, mô-đun chấm cả hai toán tử trên một tuyến rồi lấy bước tốt nhất trong các
bước được nhận, thay vì lấy bước đầu tiên.

\section{Tiêu chí chấp nhận}

\subsection{Theo trội Pareto}

Nhận bước đi khi $\mathbf{g}(\pi') \prec_\varepsilon \mathbf{g}(\pi^*)$ theo công thức~\ref{eq:dom-eps},
với $\varepsilon = 10^{-10}$. Mọi bước được nhận đều là cải thiện thật trên cả ba mục tiêu.

\subsection{Theo tổng có trọng số (mặc định)}

Nhận khi $S(\mathbf{g}(\pi')) < S(\mathbf{g}(\pi^*)) - \varepsilon$ với $S$ ở công thức~\ref{eq:score}.
Hai lựa chọn quyết định tính đúng đắn của tiêu chí này.

\paragraph{Trọng số bốc một lần cho mỗi lần gọi.} Vectơ $\mathbf{w}$ được bốc từ phân phối Dirichlet
$(1,1,1)$, tức phân bố đều trên tam giác $w_1 + w_2 + w_3 = 1$, đúng một lần ở đầu lần gọi, rồi dùng
chung cho nghiệm ban đầu và mọi ứng viên. Nếu bốc lại trọng số cho từng bước, hai bước liên tiếp có thể
được so bằng hai thước đo khác nhau và quá trình không còn đơn điệu theo bất kỳ thước đo nào. Tuỳ chọn
\code{equal} dùng $\mathbf{w} = (1/3, 1/3, 1/3)$ và không bốc số ngẫu nhiên.

\paragraph{Hệ số chuẩn hoá cố định.} Vectơ $\mathbf{s}$ được tính một lần mỗi lần chạy NSGA-II:
$s_1, s_2$ là trung bình $f_1$, $f_2$ của ba hoán vị tất định (thứ tự số hiệu, $\pi_{NN}$, $\pi_{TW}$)
chấm bằng bộ giải mã gốc; $s_3 = 1$ vì $1 - f_3$ đã nằm trong $[0,1]$. Với C102-25,
$\mathbf{s} = (3606{,}42;\ 350{,}56;\ 1)$. Hệ số này không phụ thuộc ứng viên đang xét. Chuẩn hoá theo
giá trị nhỏ nhất và lớn nhất của chính hai lời giải đang so sẽ làm thước đo thay đổi theo từng cặp,
nên không được dùng.

Một bước được nhận theo tiêu chí này có thể làm xấu một mục tiêu riêng lẻ, miễn tổng có trọng số giảm.
Vì vậy mọi bước được nhận được ghi lại kèm vectơ mục tiêu trước và sau, cùng tên tiêu chí đã dùng
(\code{accepted\_by}), để có thể kiểm tra lại.

\section{Lệch biên tuyến}
\label{sec:drift}

\subsection{Hiện tượng}

Bộ giải mã cắt hoán vị thành tuyến dựa trên tải và giờ đến tích luỹ dọc theo hoán vị. Khi đổi thứ tự
khách trong một tuyến, tải hoặc giờ đến ở cuối tuyến thay đổi, nên khách đầu của tuyến kế tiếp có thể
được ghép vào tuyến trước, hoặc ngược lại. Bộ tuyến mà bộ giải mã tạo ra từ hoán vị ứng viên khi đó khác
bộ tuyến đã đề xuất. Nếu tính lợi ích của bước đi trên bộ tuyến đề xuất thì con số đó không phải độ thích
nghi thật của cá thể.

\subsection{Phản ví dụ}

Xét ba khách nằm trên một đường thẳng, tải trọng xe $Q = 10$, nhu cầu lần lượt $6, 5, 3$ và cắt tuyến
chỉ theo tải. Hoán vị $(1,2,3)$ được cắt thành $[1] \,|\, [2,3]$ vì $6 + 5 > 10$. Đảo tuyến thứ hai bằng
2-opt cho bộ tuyến đề xuất $[1] \,|\, [3,2]$, nối lại thành hoán vị $(1,3,2)$. Bộ giải mã cắt lại
$(1,3,2)$ thành $[1,3] \,|\, [2]$ vì $6 + 3 \le 10$: khách 3 đã lọt sang tuyến đầu.

\begin{figure}[H]
\centering
\begin{tikzpicture}[
  box/.style={draw, minimum width=0.8cm, minimum height=0.7cm, font=\small},
  r1/.style={box, fill=blue1!20}, r2/.style={box, fill=orange1!25},
  lab/.style={font=\footnotesize, align=center}]
  \node[lab, anchor=east] at (-0.5,0) {hoán vị $(1,2,3)$\\ được cắt thành};
  \node[r1] at (0,0) {1}; \node[r2] at (1.1,0) {2}; \node[r2] at (1.9,0) {3};
  \draw[-{Stealth}, thick] (2.6,0) -- (3.6,0) node[midway, above, lab] {2-opt\\ tuyến 2};
  \node[r1] at (4.2,0) {1}; \node[r2] at (5.3,0) {3}; \node[r2] at (6.1,0) {2};
  \node[lab] at (5.2,-0.8) {bộ tuyến đề xuất};
  \draw[-{Stealth}, thick] (6.8,0) -- (7.8,0) node[midway, above, lab] {nối và\\ cắt lại};
  \node[r1] at (8.4,0) {1}; \node[r1] at (9.2,0) {3}; \node[r2] at (10.3,0) {2};
  \node[lab] at (9.3,-0.8) {bộ tuyến thật};
\end{tikzpicture}
\caption{Phản ví dụ lệch biên tuyến với $Q = 10$, nhu cầu $(6,5,3)$; màu thể hiện tuyến}
\label{fig:drift}
\end{figure}

\subsection{Cách xử lý}

Sau khi chấm ứng viên, mô-đun so bộ tuyến bộ giải mã trả về với bộ tuyến đề xuất, so theo đúng thứ tự
từng phần tử chứ không chỉ so tập khách. Có hai chính sách:
\begin{itemize}
  \item \code{reject} (mặc định): bộ tuyến khác thì loại bước đi. Mọi bước được nhận khi đó đúng là bước
        trong một tuyến, và bộ tuyến của hoán vị kết quả đúng bằng bộ tuyến đã đề xuất.
  \item \code{evaluate\_actual}: vẫn xét bước đi, nhưng so sánh bằng kết quả \emph{thật} của bộ giải mã
        cho hoán vị ứng viên. Không có độ thích nghi giả, vì con số dùng để so là con số NSGA-II sẽ dùng.
        Đổi lại, bước đi lúc này không còn thuần tuý trong tuyến. Sau khi nhận một bước như vậy, lượt
        duyệt dừng lại và lượt sau bắt đầu trên bộ tuyến mới.
\end{itemize}
Với cách cắt \code{capacity\_time} của nhóm, lệch biên rất thường gặp vì chỗ cắt phụ thuộc cả giờ đến:
trong thực nghiệm (mục~\ref{sec:ls-stats}), phần lớn ứng viên bị lệch biên. Đây là lý do chính sách
\code{evaluate\_actual} được đưa vào thử nghiệm.

\subsection{Khi bật cắt tuyến bằng quy hoạch động}

Khi \code{dp\_split} của TV4 bật, bộ đánh giá cắt tuyến tối ưu bằng thuật toán split của Prins~\cite{prins2004} theo vectơ trọng số riêng của từng cá
thể. Bộ cắt này được phép đổi chỗ cắt, nên mô-đun không ép bộ tuyến phải giữ nguyên. Thay vào đó, mọi
lần chấm trong một lần tìm kiếm cục bộ dùng chung đúng vectơ trọng số của con (truyền qua
\code{metadata["dp\_weights"]}); mô-đun không bốc lại trọng số DP. Bộ tuyến trả về là bộ tuyến DP thật.

\section{Ràng buộc số xe}

Bộ giải mã chỉ báo số xe có vượt $K$ hay không qua \code{fleet\_ok}, không chặn. Tìm kiếm cục bộ áp dụng
chính sách: nếu nghiệm hiện tại không vượt đội xe thì ứng viên vượt đội xe bị loại; nếu nghiệm hiện tại
đã vượt thì chỉ nhận ứng viên không dùng nhiều xe hơn. Như vậy tìm kiếm cục bộ không bao giờ biến một
nghiệm hợp lệ thành không hợp lệ, cũng không làm vi phạm nặng hơn. Mô-đun không sửa bộ đánh giá chung
hay cách NSGA-II xử lý ràng buộc.

\section{Ngân sách, bộ nhớ đệm và bộ đếm}

Mỗi lần gọi có ba trần cứng: số ứng viên được chấm (\code{max\_candidate\_evals}, tính cả ứng viên bị
loại), số bước được nhận (\code{max\_accepted\_moves}) và số lượt duyệt (\code{max\_passes}). Không có
trần theo thời gian thực, vì trần thời gian làm kết quả phụ thuộc tốc độ máy và mất tính tái lập. Do đó
một lần gọi gọi bộ đánh giá nhiều nhất $\code{max\_candidate\_evals} + 1$ lần.

Trong một lần gọi, kết quả chấm được lưu đệm theo hoán vị. Một ứng viên bốc trùng (thường gặp ở tuyến
ngắn) vẫn bị tính vào ngân sách nhưng không gọi lại bộ đánh giá. Bộ nhớ đệm chỉ sống trong một lần gọi
nên không thể lẫn kết quả giữa hai con có trọng số DP khác nhau.

Các bộ đếm cộng dồn qua một lần chạy: số lần gọi (\code{n\_ls\_calls}), số lần thực sự chạy
(\code{n\_ls\_triggered}), số ứng viên (\code{n\_ls\_candidate\_evals}), số bước được nhận, số bước bị loại
theo từng lý do (lệch biên, đội xe, không cải thiện), số yêu cầu chấm và số lần thật sự gọi bộ đánh giá
(\code{n\_real\_model\_evals}), tổng thời gian. Bộ đếm chỉ để báo cáo, không ảnh hưởng tới thuật toán.

\section{Ví dụ một lần chạy thật}

Bảng~\ref{tab:ls-run} ghi một lần tìm kiếm cục bộ trên một hoán vị ngẫu nhiên (seed 3) của C102-25, với
tiêu chí tổng có trọng số, ngân sách 40 ứng viên, tối đa 5 bước, 2 lượt. Bộ sinh của lần gọi dùng
seed 7, nên hai chính sách bốc cùng vectơ trọng số $\mathbf{w} = (0{,}307;\ 0{,}445;\ 0{,}247)$.

\begin{table}[H]
\centering
\caption{Một lần tìm kiếm cục bộ trên C102-25 với hai chính sách lệch biên}
\label{tab:ls-run}
\begin{tabular}{@{}lrr@{}}
\toprule
 & \code{reject} & \code{evaluate\_actual} \\
\midrule
$(f_1, f_2, f_3)$ ban đầu & \multicolumn{2}{c}{$(5243{,}9;\ 514{,}4;\ 0{,}3600)$} \\
$(f_1, f_2, f_3)$ sau tìm kiếm & $(5227{,}3;\ 510{,}7;\ 0{,}3600)$ & $(4799{,}0;\ 478{,}8;\ 0{,}3705)$ \\
Điểm $S$ trước $\to$ sau & $1{,}2588 \to 1{,}2528$ & $1{,}2588 \to 1{,}1731$ \\
Số ứng viên đã chấm & 30 & 20 \\
Số ứng viên bị lệch biên & 24 & 17 \\
Số ứng viên bị loại do lệch biên & 24 & 0 \\
Số bước được nhận & 2 & 2 \\
\bottomrule
\end{tabular}
\end{table}

Với \code{reject}, 24 trên 30 ứng viên bị loại vì lệch biên, hai bước được nhận chỉ giảm chi phí
khoảng 17 CNY. Với \code{evaluate\_actual}, hai bước 2-opt trên cùng một tuyến kéo theo việc cắt lại
tuyến, giảm chi phí 445 CNY, giảm khí thải 35,6 kg và tăng nhẹ mức hài lòng. Ví dụ này minh hoạ lý do
cả hai chính sách được đưa vào thực nghiệm.

\section{Giao diện lập trình}

Mô-đun nằm ở \code{src/improvements/local\_search.py}:
\begin{Verbatim}
two_opt(route, i, j) -> list[int]
or_opt(route, start, length, insert_at) -> list[int]
dominates(a, b, eps=0.0) -> bool
improve(perm, inst, rng, params, *, decode_and_evaluate, metadata=None,
        counters=None, scales=None, require_same_routes=True,
        initial_detail=None) -> LocalSearchResult
\end{Verbatim}
\code{LocalSearchResult} gồm hoán vị kết quả, kết quả đánh giá chính xác, số ứng viên, số bước được
nhận, số lượt, vectơ trọng số và lịch sử các bước được nhận. Lớp \code{LocalSearchHook} trong
\code{tv2\_hooks.py} bọc hàm này cho NSGA-II: mỗi lần gọi bốc đúng một số để quyết định có chạy hay
không, tính $\mathbf{s}$ một lần khi khởi tạo và giữ bộ đếm.
